\chapter{ماذا يقول الدوبامين حقًا}
% full version: chapters/07_dofaalt.tex

في الفصل السابق كان الدوبامين بسيطًا: عدد واحد، ”أفضل أو أسوأ مما كان متوقعًا“. هذا تقريب أول، ويعمل جيدًا. لكن تجارب السنوات الأخيرة بيّنت أن ما يجري في هذا ”السلك“ أكثر من عدد واحد. هذا الفصل عن المواضع التي تتعقد فيها الصورة ولا يزال العلماء يتجادلون فيها.

\section*{المتفائلون والمتشائمون}

عصبونات الدوبامين ليست متماثلة. بعضها يستجيب للمفاجآت السارة أقوى من المزعجة: هؤلاء المتفائلون. وبعضها بالعكس: المتشائمون. إذا تعلم كل عصبون بقاعدته الخاصة، عرف المتفائل كيف ستكون المكافأة في أحسن الأحوال، والمتشائم كيف ستكون في أسوئها. ومعًا لا يسجلون توقعًا متوسطًا واحدًا، بل انتشار النتائج الممكنة كله، كوكيل مراهنات لا يعرف المرشح الأوفر حظًا فحسب، بل فرص كل متسابق. توافق هذه الفكرة التجربة؛ أما لماذا يحتاج الدماغ إلى الانتشار كله فلا يزال فرضية.

\pencilfig{7}{\pencilart{7}}{كل عصبون دوبامين يمسك نقطته الخاصة على منحنى المكافآت الممكنة: ومعًا يتذكرون المنحنى كله.}

\section*{لا ”أفضل“ فحسب، بل ”مهم“ أيضًا}

بعض عصبونات الدوبامين يطلق رشقة حتى عند أحداث مزعجة لكنها مهمة: صوت عالٍ، خطر. يبدو أن شبكة الدوبامين لا تحمل اتجاه المفاجأة (أفضل أم أسوأ) فحسب، بل أهميتها أيضًا: إشارة ”انتبه“.

\section*{إشارتان على سلك واحد}

الرشقات السريعة تعلّم. أما المستوى المتوسط البطيء للدوبامين، الذي يتغير خلال دقائق، فيبدو أنه يقول شيئًا آخر: كم هي البيئة غنية الآن. إذا كثرت المكافآت حولك، غلا الوقت، فيتصرف الحيوان أسرع. وتؤكد التجارب أن مستوى الدوبامين مرتبط بحيوية الأفعال. المهندس سيسمي هذا فصل الإشارات بحسب التردد: الترددات العالية تحمل شيئًا، والمنخفضة شيئًا آخر.

\section*{التصاعد عند الهدف}

حين يقترب الحيوان من المكافأة، كثيرًا ما لا يطلق الدوبامين رشقة، بل يرتفع ارتفاعًا سلسًا. تفسير أول: هذا هو الترقب نفسه، إشارة ”الهدف قريب، شدّ حيلك“. وتفسير ثانٍ يحتفظ بفكرة المفاجأة: من بعيد يقدّر الحيوان وضعه تقديرًا ضبابيًا، ومع الاقتراب يزداد دقة، وكل تدقيق مفاجأة سارة صغيرة. وتجربة أنيقة في ممر افتراضي، نُقل فيها الحيوان فجأة إلى موضع أقرب من الهدف، تشهد للتفسير الثاني: استجاب الدوبامين برشقة، كما يليق بالمفاجأة.

\section*{نظرة إلى الوراء}

وثمة بديل جريء: الدوبامين لا يجيب عن سؤال ”ماذا سيحدث؟“، بل عن سؤال ”ما الذي تسبب في المكافأة؟“، أي ينظر إلى الماضي لا إلى المستقبل. والتجارب التي تتنبأ فيها النظريتان بنتائج مختلفة لا تزال متعارضة. إنه جدل علمي حي.

\section*{متجه بدل عدد}

وأخيرًا، يستجيب الدوبامين أيضًا لتغير \emph{نكهة} المكافأة حين لا تتغير قيمتها، ويساعد على التعلم بلا أي مكافأة. يبدو أنه ليس خطأ واحدًا، بل حزمة كاملة من الأخطاء: خطأ لكل سمة مما حدث.

خلاصة الفصل: معظم النظريات الجديدة لا تلغي الصورة البسيطة، بل توسعها. الدوبامين ليس سلكًا واحدًا، بل حزمة صغيرة من عدة أسلاك. وحدها ”النظرة إلى الوراء“ تجادلها جدالًا جادًا.

\fullversion{7}
